\input{sol_pre.tex}
\input{q_2022-12-04.tex}
\begin{document}
\input{sol_head.tex}

%================= 1
\Savol{1}
Son eng katta bölişi uçun yuzlar xonasiga eng katta raqam: $a=9$. Raqamlar turli, demak $b=8$, $c=7$.
\[100a+10b+c=987.\]
\Javob{B) 987}

%================= 2
\Savol{2}
\[2022-\frac1{2022}=\frac{2022^2-1}{2022}=\frac{2021\cdot2023}{2022}.\]
\[\frac{2021\cdot2023}{2022}\cdot\frac{2022}{2021}\cdot\frac1{2023}=1.\]
\Javob{A) 1}

%================= 3
\Savol{3}
Ikki xonali sonlar: $10,\dots,99$ (90 ta), yiğindisi $\dfrac{(10+99)\cdot 90}{2}=4905$.\\
2 ga karralilar: $10,12,\dots,98$ — 45 ta, $S_2=\dfrac{(10+98)\cdot45}{2}=2430$.\\
3 ga karralilar: $12,15,\dots,99$ — 30 ta, $S_3=\dfrac{(12+99)\cdot30}{2}=1665$.\\
6 ga karralilar (ikki marta ayrilgan): $12,\dots,96$ — 15 ta, $S_6=\dfrac{(12+96)\cdot15}{2}=810$.
\[S=4905-2430-1665+810=1620.\]
\Javob{A) 1620}

%================= 4
\Savol{4}
$AB=S$, uçraşgunça vaqt $t$. Birinçi maşina $B$ gaça borib qaytgan, demak ikkalasining yöllari yiğindisi $2S$:
\[v_1t+v_2t=2S\ \Rightarrow\ t=\frac{2S}{v_1+v_2},\qquad v_2t=\frac{2v_2}{v_1+v_2}\,S.\]
\Javob{A) $\dfrac{2v_2}{v_1+v_2}\cdot100\%$}

%================= 5
\Savol{5}
Bir soatda bajariladigan iş (hovuz $=1$):
\[\frac15+\frac13-\frac12=\frac{6+10-15}{30}=\frac1{30}\ \Rightarrow\ t=30 \text{ soat}.\]
\Javob{B) 30}

%================= 6
\Savol{6}
$38^2-23^2=(38-23)(38+23)=15\cdot61$. U holda
\[\sqrt{\frac{244}{15\cdot15\cdot61}}=\sqrt{\frac{4\cdot61}{225\cdot61}}=\frac{2}{15},\qquad
\sqrt[3]{\frac{12}{25}\cdot\frac{2}{15}}=\sqrt[3]{\frac{8}{125}}=\frac25=0{,}4.\]
\Javob{D) 0,4}

%================= 7
\Savol{7}
Geometrik progressiyada $b_2^2=b_1b_3$:
\[(a+4)^2=(a-5)(5a+8)\ \Rightarrow\ 4a^2-25a-56=0\ \Rightarrow\ a=8\ \text{yoki}\ a=-\tfrac74.\]
$a=8$: $3,\,12,\,48$ — işora almaşmaydi (mos emas).\\
$a=-\frac74$: $b_1=-\frac{27}4,\ b_2=\frac94,\ b_3=-\frac34$, $q=-\frac13$. Demak $b_4=b_3q=\frac14$.
\Javob{A) $\frac14$}

%================= 8
\Savol{8}
\[\frac{1}{(x+y)^2}\cdot\frac{x^2+y^2}{x^2y^2}+\frac{2}{(x+y)^3}\cdot\frac{x+y}{xy}
=\frac{x^2+y^2+2xy}{x^2y^2(x+y)^2}=\frac{(x+y)^2}{x^2y^2(x+y)^2}=\frac1{x^2y^2}.\]
\Javob{B) $\frac1{x^2y^2}$}

%================= 9
\Savol{9}
$D=b^2-4ac=25-32=-7<0$ — haqiqiy ildiz yöq.
\Javob{A) Haqiqiy ildizga ega emas}

%================= 10
\Savol{10}
$\sin\frac{7\pi}6=\sin\!\left(\pi+\frac\pi6\right)=-\frac12$, \quad $\cos\frac{7\pi}4=\cos\!\left(2\pi-\frac\pi4\right)=\frac{\sqrt2}2$.
\[2\cdot\left(-\tfrac12\right)+\left(\tfrac{\sqrt2}{2}\right)^2=-1+\tfrac12=-\tfrac12.\]
\Javob{B) $-\frac12$}

%================= 11
\Savol{11}
$\sqrt{u^2}=|u|$. $a>5$ da: $|3-a|=a-3$, $|a-5|=a-5$.
\[(a-3)-(a-5)=2.\]
\Javob{C) 2}

%================= 12
\Savol{12}
$b=10c$, $a=4b=40c$. U holda
\[\frac{a^2+b^2+c^2}{ac}=\frac{(1600+100+1)c^2}{40c^2}=\frac{1701}{40}=42\tfrac{21}{40},\qquad \frac ac=40.\]
Yiğindi: $42\frac{21}{40}+40=82\frac{21}{40}$.
\Javob{B) $82\frac{21}{40}$}

%================= 13
\Savol{13}
Surat: $\cos3\alpha+\cos5\alpha=2\cos4\alpha\cos\alpha$, demak
\[2\cos\alpha+2\cos4\alpha\cos\alpha=2\cos\alpha(1+\cos4\alpha)=2\cos\alpha\cdot2\cos^22\alpha.\]
Maxraj: $\cos3\alpha=\cos2\alpha\cos\alpha-\sin2\alpha\sin\alpha$, şuning uçun
\[\cos3\alpha+\sin\alpha\sin2\alpha=\cos2\alpha\cos\alpha.\]
\[\frac{4\cos\alpha\cos^22\alpha}{\cos2\alpha\cos\alpha}=4\cos2\alpha.\]
\Javob{B) $4\cos2\alpha$}

%================= 14
\Savol{14}
Tenglamani $4^x$ ga bölamiz va $t=\left(\frac52\right)^x>0$ deymiz ($\frac{25^x}{4^x}=t^2$, $\frac{10^x}{4^x}=t$):
\[2t^2-5t+3=0\ \Rightarrow\ t=1\ \text{yoki}\ t=\tfrac32.\]
Ikkalasi ham musbat, har biri bittadan $x$ beradi ($x=0$ va $x=\log_{2{,}5}1{,}5$). Demak 2 ta ildiz.
\Javob{C) 2 ta}

%================= 15
\Savol{15}
Toq funksiya: $f(-x)=-f(x)$. A) da surat $\sin x+x^3$ — toq, maxraj $\cos x-1$ — juft:
\[f(-x)=\frac{-\sin x-x^3}{\cos x-1}=-f(x).\]
B), C), D) da surat yoki maxraj na juft, na toq (masalan, $\sin x-1$), şuning uçun ular toq emas.
\Javob{A)}

%================= 16
\Savol{16}
$t=\sqrt{x^2+4x+4}=|x+2|\ge0$ deymiz. $x^2+4x+1=t^2-3$:
\[t^2-3=2t\ \Rightarrow\ t^2-2t-3=0\ \Rightarrow\ t=3\ (t=-1 \text{ mos emas}).\]
$|x+2|=3\Rightarrow x=1$ yoki $x=-5$. Yiğindi: $1+(-5)=-4$.
\Javob{A) $-4$}

%================= 17
\Savol{17}
\[-5\le|2x-1|-7\le5\ \Rightarrow\ 2\le|2x-1|\le12.\]
1) $2\le2x-1\le12\Rightarrow 1{,}5\le x\le6{,}5$: butunlar $2,3,4,5,6$ — 5 ta.\\
2) $-12\le2x-1\le-2\Rightarrow -5{,}5\le x\le-0{,}5$: butunlar $-5,\dots,-1$ — 5 ta.\\
Jami: $5+5=10$.
\Javob{C) 10 ta}

%================= 18
\Savol{18}
Arifmetik ildiz har doim $\ge0$, şuning uçun tengsizlik aniqlaniş sohasining hamma nuqtasida bajariladi:
\[4-x^2\ge0\ \Rightarrow\ -2\le x\le2.\]
\Javob{A) $[-2;\,2]$}

%================= 19
\Savol{19}
\[32^3\cdot8^{12}\cdot125^{17}=2^{15}\cdot2^{36}\cdot5^{51}=2^{51}\cdot5^{51}=10^{51}.\]
$10^{51}$ — 1 dan keyin 51 ta nol, ya'ni 52 xonali son.
\Javob{A) 52}

%================= 20
\Savol{20}
Logarifm asosi: $0<\cos2x\ne1$, ifoda: $\sin2x>0$. $1=\log_{\cos2x}\cos2x$ va asos $0<\cos2x<1$ bölgani uçun funksiya kamayuvçi — işora almaşadi:
\[\sin2x\ge\cos2x>0.\]
$\sin2x>0,\ \cos2x>0$ — $2x$ I çorakda; ularda $\sin\ge\cos$ bölişi uçun $\frac\pi4\le 2x<\frac\pi2$ ($2x=\frac\pi2$ da $\cos2x=0$ — mumkin emas):
\[\frac\pi8+\pi k\le x<\frac\pi4+\pi k,\quad k\in\mathbb Z.\]
\Javob{C)}

%================= 21
\Savol{21}
$g^{-1}(g(x))=x$, demak $f(x+1)=x^2+5x+6$. $x+1=-1\Rightarrow x=-2$:
\[f(-1)=4-10+6=0.\]
\Javob{D) 0}

%================= 22
\Savol{22}
$x=-\frac\pi6$ da modul içidagi ifoda: $\cos\!\left(-\frac\pi3\right)-\frac\pi6=0{,}5-0{,}5236<0$. Demak bu nuqta atrofida
\[f(x)=\sin^25x+\cos2x+x,\qquad f'(x)=5\sin10x-2\sin2x+1.\]
\[f'\!\left(-\tfrac\pi6\right)=5\sin\!\left(-\tfrac{5\pi}3\right)-2\sin\!\left(-\tfrac\pi3\right)+1=\tfrac{5\sqrt3}2+\sqrt3+1=\tfrac{7\sqrt3+2}2.\]
\Javob{A) $\frac{7\sqrt3+2}{2}$}

%================= 23
\Savol{23}
\fig{\begin{tikzpicture}[scale=1.05]
\coordinate (A) at (0,0); \coordinate (C) at (5,0); \coordinate (B) at (2.75,2.75);
\coordinate (F) at (1.25,0); \coordinate (E) at ($(C)!0.75!(B)$);
\fill[aux!12] (F)--(E)--(C)--cycle;
\draw[thick] (A)--(B)--(C)--cycle; \draw[axs] (F)--(E); \draw[axs] (F)--(C);
\node[left] at (A){$A$}; \node[above] at (B){$B$}; \node[right] at (C){$C$};
\node[below] at (F){$F$}; \node[above right] at (E){$E$};
\node[above left] at ($(A)!0.5!(B)$){$16$}; \node[above left,aux] at ($(F)!0.5!(E)$){$12$};
\node[above,aux] at ($(F)!0.55!(C)$){$x$};
\draw[|<->|] (0,-0.6)--node[lbl]{$20$}(5,-0.6);
\end{tikzpicture}}
$EF\parallel AB$ bölgani uçun $\triangle FEC\sim\triangle ABC$ (\textcolor{aux}{qizil} uçburçak):
\[\frac{FC}{AC}=\frac{EF}{AB}\ \Rightarrow\ \frac{x}{20}=\frac{12}{16}\ \Rightarrow\ x=15.\]
\Javob{B) 15}

%================= 24
\Savol{24}
\fig{\begin{tikzpicture}[scale=1.1]
\draw[fill=msblue!6] (-2.6,-1.2)--(3.2,-1.2)--(4.2,1.3)--(-1.6,1.3)--cycle;
\coordinate (H) at (0.3,0.6); \coordinate (B) at (0.3,3.3); \coordinate (A) at (-1.7,-0.4); \coordinate (C) at (2.5,-0.5);
\draw[axs] (A)--(H)--(C); \draw[ax] (A)--(C);
\draw[thick] (A)--(B)--(C); \draw[thick] (B)--(H);
\draw[aux] ($(H)!0.2!(A)$)--($(H)!0.2!(A)+(H)!0.17!(C)-(H)$)--($(H)!0.17!(C)$);
\draw ($(H)+(0,0.3)$)--($(H)+(0.25,0.22)$)--($(H)+(0.25,-0.08)$);
\foreach \p in {A,B,C,H} \fill (\p) circle (1.3pt);
\node[left] at (A){$A$}; \node[above] at (B){$B$}; \node[right] at (C){$C$}; \node at ($(H)+(-0.3,0.2)$){$H$};
\node[left] at ($(B)!0.55!(H)$){$a$}; \node[aux] at ($(H)!0.5!(C)+(-0.15,-0.25)$){$a$};
\node[above left] at ($(A)!0.5!(B)$){$2a$}; \node[aux,below] at ($(A)!0.5!(C)$){$2a$};
\node[above right] at ($(B)!0.5!(C)$){$a\sqrt2$};
\end{tikzpicture}}
$BH=HC=a$, $AB=2a$. $BH$ tekislikka perpendikulyar, demak $\triangle BHA$ va $\triangle BHC$ töğri burçakli:
\[AH=\sqrt{4a^2-a^2}=a\sqrt3,\qquad BC=\sqrt{a^2+a^2}=a\sqrt2.\]
$\angle AHC=90^\circ$: \ $AC=\sqrt{3a^2+a^2}=2a$. \ $\triangle ABC$ da kosinuslar teoremasi:
\[\cos\angle BAC=\frac{AB^2+AC^2-BC^2}{2\cdot AB\cdot AC}=\frac{4a^2+4a^2-2a^2}{2\cdot2a\cdot2a}=\frac34.\]
\Javob{A) $\frac34$}

%================= 25
\Savol{25}
\fig{\begin{tikzpicture}[scale=0.55]
\draw[thick] (0,0) rectangle (4,4); \draw[thick,fill=msblue!8] (2,2) circle (2);
\fill (2,2) circle (2pt); \draw[axs] (2,2)--(4,2) node[midway,above,aux]{$r=2$};
\node[below] at (2,0){$4$};
\end{tikzpicture}}
Içki çizilgan doira diametri kvadrat tomoniga teng: $2r=4\Rightarrow r=2$. \quad $S=\pi r^2=4\pi$.
\Javob{D) $4\pi$}

%================= 26
\Savol{26}
\fig{\begin{tikzpicture}[scale=1.15]
\coordinate (A) at (0,0); \coordinate (B) at (0,2.449); \coordinate (C) at (4.243,0); \coordinate (D) at (1.414,0);
\draw[thick] (A)--(B)--(C)--cycle; \draw[thick] (B)--(D);
\draw (0.22,0)--(0.22,0.22)--(0,0.22);
\pic[draw,angle radius=0.55cm,"\scriptsize$30^\circ$",angle eccentricity=1.75] {angle=A--B--D};
\pic[draw,angle radius=0.75cm,"\scriptsize$30^\circ$",angle eccentricity=1.45] {angle=D--B--C};
\pic[draw,angle radius=0.8cm,"\scriptsize$30^\circ$",angle eccentricity=1.5] {angle=B--C--A};
\node[left] at (A){$A$}; \node[above] at (B){$B$}; \node[right] at (C){$C$}; \node[below] at (D){$D$};
\node[lbl] at ($(B)!0.5!(D)+(0.5,0)$){$2\sqrt2$};
\end{tikzpicture}}
$\angle B=90^\circ-30^\circ=60^\circ$, bissektrisa uni ikkiga böladi: $\angle ABD=30^\circ$.\\
$\triangle ABD$ ($\angle A=90^\circ$): \ $AB=BD\cos30^\circ=2\sqrt2\cdot\frac{\sqrt3}2=\sqrt6$.\\
$\triangle ABC$: \ $AC=AB\cdot\tg60^\circ=\sqrt6\cdot\sqrt3=3\sqrt2$.
\[S=\frac12\cdot AB\cdot AC=\frac12\cdot\sqrt6\cdot3\sqrt2=3\sqrt3.\]
\Javob{D) $3\sqrt3$}

%================= 27
\Savol{27}
$n$ burçakli köpburçak diagonallari soni $\dfrac{n(n-3)}2$:
\[\frac{n(n-3)}2=6n\ \Rightarrow\ n-3=12\ \Rightarrow\ n=15.\]
Içki burçaklar yiğindisi: $180^\circ(n-2)=180^\circ\cdot13=2340^\circ$.
\Javob{C) $2340^\circ$}

%================= 28
\Savol{28}
\fig{\begin{tikzpicture}[scale=0.85]
\coordinate (A) at (0,0); \coordinate (B) at (4,0); \coordinate (D) at (2.828,2.828); \coordinate (C) at (6.828,2.828);
\coordinate (O) at ($(A)!0.5!(C)$);
\draw[thick] (A)--(B)--(C)--(D)--cycle; \draw (A)--(C); \draw (B)--(D);
\foreach \X/\Y in {A/B,B/C,C/D,D/A}{\coordinate (P) at ($(\X)!(O)!(\Y)$); \draw[axs] (O)--(P);}
\fill (O) circle (1.5pt);
\node[below left] at (A){$A$}; \node[below right] at (B){$B$}; \node[above right] at (C){$C$}; \node[above left] at (D){$D$};
\node at ($(O)+(-0.4,0.13)$){$O$};
\node[aux] at ($(O)+(-0.25,-0.9)$){$5$};
\pic[draw,angle radius=0.7cm,"\scriptsize$45^\circ$",angle eccentricity=1.5] {angle=B--A--D};
\end{tikzpicture}}
Romb tomoni $a$. Katta diagonal: $d_1^2=a^2+a^2-2a^2\cos135^\circ=a^2(2+\sqrt2)$, ya'ni
\[d_1=a\sqrt{2+\sqrt2}=10\sqrt{4+2\sqrt2}=10\sqrt2\cdot\sqrt{2+\sqrt2}\ \Rightarrow\ a=10\sqrt2.\]
Romb balandligi $h=a\sin45^\circ=10$. $O$ dan har bir tomongaça masofa $\frac h2=5$ (\textcolor{aux}{qizil} perpendikulyarlar). Törtala tomongaça: $4\cdot5=20$.
\Javob{B) 20}

%================= 29
\Savol{29}
$2x^2|2x|=4x^2|x|$ — juft funksiya. $|x|$ ni oçamiz:
\[\int_{-3}^{4}4x^2|x|\,dx=\int_{-3}^{0}(-4x^3)\,dx+\int_{0}^{4}4x^3\,dx=x^4\Big|_{0}^{3}\!+x^4\Big|_{0}^{4}=81+256=337.\]
\Javob{C) 337}

%================= 30
\Savol{30}
\fig{\begin{tikzpicture}[scale=0.9]
\coordinate (A) at (0,0); \coordinate (B) at (2.4,0); \coordinate (E) at (4.8,0); \coordinate (C) at (1.4,2.2);
\fill[aux!10] (B)--(E)--(C)--cycle;
\draw[thick] (A)--(B)--(C)--cycle; \draw[axs] (B)--(E)--(C);
\draw[ax] (C)--(1.4,0); \draw[ax] (1.4,0.18)--(1.58,0.18)--(1.58,0);
\node[below] at (A){$A$}; \node[below] at (B){$B$}; \node[below] at (E){$E$}; \node[above] at (C){$C$};
\node[aux,lbl] at (1.62,1){$h$};
\end{tikzpicture}}
$\overrightarrow{BE}=\overrightarrow{AB}$ — $E$ nuqta $AB$ nurida, $BE=AB$, ya'ni $AE=2AB$.\\
$\triangle AEC$ va $\triangle ABC$ ning $C$ dan tuşirilgan balandligi $h$ umumiy, asosi 2 marta katta:
\[S_{AEC}=2S_{ABC}=120.\]
\Javob{A) 120}

%================= 31
\Savol{31}
\fig{\begin{tikzpicture}[scale=0.62]
\fill[aux!15] (0,0) rectangle (10,5.2);
\fill[white] (3.7,2.6) circle (2) (6.3,2.6) circle (2);
\begin{scope}\clip (3.7,2.6) circle (2); \fill[aux!15] (6.3,2.6) circle (2);\end{scope}
\draw (0,0) rectangle (10,5.2); \draw (3.7,2.6) circle (2) (6.3,2.6) circle (2);
\node at (2,4.6){$A$}; \node at (8,4.6){$B$}; \node at (9.6,0.4){$U$};
\foreach \x/\y/\t in {0.9/4.4/10, 0.8/3.2/11, 1.3/2.2/4, 0.8/1.0/9}{\node[aux] at (\x,\y){\bfseries\t};}
\foreach \x/\y/\t in {3.4/3.8/12, 2.9/2.5/1, 3.7/1.2/8, 6.6/3.6/6, 7.2/2.5/14, 6.8/1.2/17}{\node at (\x,\y){\t};}
\foreach \x/\y/\t in {5/3.3/3, 5/2.6/5, 5/1.9/7}{\node[aux] at (\x,\y){\bfseries\t};}
\end{tikzpicture}}
$A\cap B=\{3;\,5;\,7\}$ — ikkala doiraning umumiy qismi.\\
$(A\cup B)'$ — ikkala doiradan taşqaridagi elementlar: $\{10;\,11;\,4;\,9\}$.\\
Birlaşma (\textcolor{aux}{qizil} sohalar): $3+4=7$ ta element.
\Javob{B) 7}

%================= 32
\Savol{32}
Otişlar boğliq emas. Tegiş $p=0{,}8$, tegmaslik $q=0{,}2$:
\[P=0{,}8^3\cdot0{,}2^2=0{,}512\cdot0{,}04=0{,}02048=2048\cdot10^{-5}=2^{11}\cdot10^{-5}.\]
\Javob{B) $2^{11}\cdot10^{-5}$}

%================= 33-35
\newpage
\Savol{33}
\fig{\begin{tikzpicture}[scale=0.5]
% 3D piramida
\begin{scope}
\coordinate (O) at (0,0); \coordinate (S) at (0,5.2);
\foreach \i in {0,...,5}{\coordinate (P\i) at ({3*cos(60*\i+10)},{1.0*sin(60*\i+10)});}
\draw[thick] (P5)--(P0)--(P1); \draw[thick] (P2)--(P3)--(P4)--(P5);
\draw[dashed] (P1)--(P2);
\foreach \i in {0,3,4,5} \draw[thick] (S)--(P\i);
\foreach \i in {1,2} \draw[dashed] (S)--(P\i);
\coordinate (M) at ($(P4)!0.5!(P5)$);
\draw[ax] (S)--(O); \draw[ax] (O)--(M); \draw[axs] (S)--(M);
\fill (O) circle (1.5pt); \fill (M) circle (1.5pt);
\node[above] at (S){$S$}; \node at ($(O)+(-0.38,0.22)$){$O$}; \node[below] at (M){$M$};
\node[below] at ($(P4)!0.5!(M)$){\small$6$};
\end{scope}
% kesim
\begin{scope}[xshift=9cm,yshift=-0.6cm]
\coordinate (M1) at (-3.118,0); \coordinate (M2) at (3.118,0); \coordinate (S) at (0,5.4); \coordinate (O) at (0,0);
\coordinate (I) at (0,1.8);
\draw[thick] (M1)--(S)--(M2)--cycle; \draw[ax] (S)--(O);
\draw[aux,thick,fill=aux!6] (I) circle (1.8);
\coordinate (T) at (0,3.6); \coordinate (K1) at (-1.697,1.2); \coordinate (K2) at (1.697,1.2);
\draw[okgreen,thick,fill=okgreen!12] (T)--(K1)--(K2)--cycle;
\draw[ax] (I)--(K2); \draw[ax] (M2)--(I);
\fill (I) circle (1.5pt); \fill (O) circle (1.5pt);
\node[above] at (S){$S$}; \node[below] at (O){$O$}; \node[below] at (M2){$M$}; \node[below] at (M1){$M'$};
\node[left] at (I){$I$};
\node[below,aux] at ($(O)!0.5!(M2)+(0,-0.05)$){\small$3\sqrt3$};
\pic[draw,angle radius=0.45cm] {angle=S--M2--M1};
\end{scope}
\end{tikzpicture}}
\textbf{Asos.} Muntazam oltiburçak tomoni $a=6$, içki çizilgan aylana radiusi (apofema) $OM=\frac{a\sqrt3}2=3\sqrt3$, yuzi
$S_{\text{asos}}=\frac{3\sqrt3}{2}a^2=54\sqrt3$.

\textbf{33.} $\triangle SOM$: $\angle SMO=60^\circ$ (yon yoq asos bilan), $H=SO=OM\cdot\tg60^\circ=3\sqrt3\cdot\sqrt3=9$.
\[V_{\text{pir}}=\frac13\cdot54\sqrt3\cdot9=162\sqrt3.\]
\textbf{34.} Şar markazi $I$ — $SO$ da, $MI$ esa $\angle SMO$ ning bissektrisasi (\textcolor{aux}{qizil} kesim): $\rho=OM\cdot\tg30^\circ=3\sqrt3\cdot\frac1{\sqrt3}=3$.
\[V_{\text{şar}}=\frac43\pi\rho^3=\frac43\cdot3\cdot27=108.\]
\textbf{35.} Radiusi $\rho$ bölgan şarga içki çizilgan konus hajmi $h$ balandlikda: asos radiusi $r^2=\rho^2-(h-\rho)^2=6h-h^2$,
\[V(h)=\frac\pi3(6h^2-h^3),\quad V'(h)=\frac\pi3(12h-3h^2)=0\Rightarrow h=4,\ r^2=8.\]
\[V_{\text{konus}}=\frac13\pi\cdot8\cdot4=\frac{32\pi}3\approx32\quad(\textcolor{okgreen}{\text{yaşil}}\ \text{konus}).\]
\Javob{33 – A; \ 34 – B; \ 35 – C}

%================= 36
\newpage
\Savol{36}
Aniqlaniş sohasi: $x\ne\pm1$. Proporsiya xossasi:
\[(6x-k)(x+1)=(3x+k)(x-1)\]
\[6x^2+6x-kx-k=3x^2-3x+kx-k\ \Rightarrow\ 3x^2+(9-2k)x=0\ \Rightarrow\ x\,(3x+9-2k)=0.\]
Ildizlar: \quad $x_1=0$ (har doim mos, çunki $0\ne\pm1$), \quad $x_2=\dfrac{2k-9}{3}$.

Tenglama \textbf{bitta} ildizga ega böladi, agar:
\begin{enumerate}\setlength{\itemsep}{2pt}
\item $x_2=x_1$: \ $\dfrac{2k-9}{3}=0\Rightarrow k=4{,}5$;
\item $x_2$ çet ildiz bölsa, ya'ni $x_2=1$ yoki $x_2=-1$:
\[\frac{2k-9}{3}=1\Rightarrow k=6;\qquad \frac{2k-9}{3}=-1\Rightarrow k=3.\]
\end{enumerate}
Tekşiruv: $k=6$ da $x(3x-3)=0\Rightarrow x=0$ yoki $x=1$ (çet) — bitta ildiz; $k=3$ da $x(3x+3)=0\Rightarrow x=0$ yoki $x=-1$ (çet) — bitta ildiz.

Demak $k\in\{3;\ 4{,}5;\ 6\}$.

a) Eng katta qiymat: $k=6$.\\
b) Yiğindi: $3+4{,}5+6=13{,}5$.
\Javob{a) 6; \quad b) 13,5}

%================= 37
\newpage
\Savol{37}
Aniqlaniş sohasi: $\sin x>0$ (u holda $1-\cos2x=2\sin^2x>0$ ham bajariladi).\\
$1-\cos2x=2\sin^2x$, demak
\[\log_2(2\sin^2x)=1+2\log_2\sin x.\]
$t=\log_2\sin x$ deymiz. $0<\sin x\le1$ bölgani uçun $t\le0$.
\[3t^2+1+2t=2\ \Rightarrow\ 3t^2+2t-1=0\ \Rightarrow\ t=-1\ \text{yoki}\ t=\tfrac13.\]
$t=\frac13>0$ — mos emas. \ $t=-1$: \ $\sin x=\frac12$:
\[x=\frac\pi6+2\pi n,\qquad x=\frac{5\pi}6+2\pi n,\quad n\in\mathbb Z.\]
\fig{\begin{tikzpicture}[scale=1.5]
\draw[->] (-1.3,0)--(1.4,0) node[right]{\small$x$}; \draw[->] (0,-1.2)--(0,1.3) node[above]{\small$y$};
\draw[thick] (0,0) circle (1);
\draw[axs] (-1.15,0.5)--(1.15,0.5); \node[aux,right] at (1.15,0.5){\small$y=\frac12$};
\draw[ax] (0,0)--(0.866,0.5); \draw[ax] (0,0)--(-0.866,0.5);
\fill[aux] (0.866,0.5) circle (1.3pt); \fill[aux] (-0.866,0.5) circle (1.3pt);
\node[above right] at (0.866,0.5){\small$\frac\pi6$}; \node[above left] at (-0.866,0.5){\small$\frac{5\pi}6$};
\end{tikzpicture}}
a) Eng kiçik musbat ildiz: $x=\dfrac\pi6$.

b) $\left(-\frac\pi2;\ \frac{5\pi}2\right)$ oraliqda: \ $\frac\pi6$, \ $\frac{5\pi}6$, \ $\frac\pi6+2\pi=\frac{13\pi}6$ \ ($\frac{13\pi}6<\frac{15\pi}6=\frac{5\pi}{2}$). \ $\frac{5\pi}6+2\pi=\frac{17\pi}6>\frac{5\pi}2$ — kirmaydi, manfiy tomonda $-\frac{7\pi}{6}<-\frac{\pi}{2}$ — kirmaydi. Jami 3 ta.
\Javob{a) $\frac\pi6$; \quad b) 3}

%================= 38
\newpage
\Savol{38}
\fig{\begin{tikzpicture}[scale=0.62]
\draw[->] (-3.2,0)--(7.4,0) node[right]{$x$}; \draw[->] (0,-2.6)--(0,7.2) node[above]{$y$};
\draw[thick,domain=-2.6:4.6,samples=60] plot (\x,{-0.5*\x*\x+\x+4});
\draw[thick,domain=-1:7] plot (\x,{-\x+6});
\draw[ax] (2,0)--(2,4); \draw[ax] (0,4)--(2,4);
\fill (2,4) circle (2pt); \fill (-2,0) circle (2pt); \fill (0,4) circle (2pt); \fill (0,6) circle (2pt); \fill (6,0) circle (2pt);
\node[below] at (2,0){$2$}; \node[below left] at (-2,0){$-2$}; \node[below] at (6,0){$6$};
\node[left] at (0,4.35){$4$}; \node[left] at (0,6){$6$};
\node[above right] at (2.05,4.05){$(2;4)$};
\node at (-2.6,2.4){$f(x)$}; \node at (5.2,1.8){$g(x)$};
\end{tikzpicture}}
\textbf{Töğri çiziq} $(0;6)$ va $(6;0)$ dan ötadi: $g(x)=-x+6$, burçak koeffitsiyenti $k=-1$.\\
$g$ parabolaga $x=2$ nuqtada urinadi: $f(2)=g(2)=4$, \ $f'(2)=k=-1$.

a) $f'(2)=-1$.

b) \textbf{Parabola}: $f(0)=4\Rightarrow c=4$. \ $f'(x)=2ax+b$.
\[\begin{cases} f(-2)=4a-2b+4=0\\ f'(2)=4a+b=-1\end{cases}\ \Rightarrow\ 3b=3,\ b=1,\ a=-\tfrac12.\]
($f(2)=-2+2+4=4$ — töğri.) \ $f(x)=-\frac{x^2}2+x+4$, \ $f'(x)=-x+1$.
\[\frac1a=-2,\ \ \frac1b=1,\ \ \frac1c=\frac14:\qquad f'(-2)-5f'(1)+2f'\!\left(\tfrac14\right)=3-5\cdot0+2\cdot\tfrac34=4{,}5.\]
\Javob{a) $-1$; \quad b) 4,5}

%================= 39
\newpage
\Savol{39}
Köphadni köpaytuvçilarga ajratamiz:
\[x^4-2x^3+2x^2-2x+1=(x^4-2x^3+x^2)+(x^2-2x+1)=x^2(x-1)^2+(x-1)^2=(x^2+1)(x-1)^2.\]
a) $x=1\pm\sqrt2$ da $(x-1)^2=2$, \ $x^2=(1\pm\sqrt2)^2=3\pm2\sqrt2$, \ $x^2+1=4\pm2\sqrt2$:
\[f(1+\sqrt2)+f(1-\sqrt2)=2(4+2\sqrt2)+2(4-2\sqrt2)=16.\]
b) $(x^2+1)(x-1)^2=0$. \ $x^2+1>0$ — ildizi yöq; \ $(x-1)^2=0\Rightarrow x=1$.\\
Turli haqiqiy ildizlar soni: 1 ta.
\Javob{a) 16; \quad b) 1}

%================= 40
\newpage
\Savol{40}
\fig{\begin{tikzpicture}[xscale=0.55,yscale=0.4]
\fill[msblue!15,domain=-2:8,samples=60] plot (\x,{0.6*\x+1.2}) -- plot[domain=8:-2] (\x,{(\x*\x-3*\x-10)/5});
\draw[->] (-3.2,0)--(9.6,0) node[right]{$x$}; \draw[->] (0,-3.4)--(0,8.4) node[above]{$y$};
\draw[thick,domain=-2.6:8.5,samples=60] plot (\x,{(\x*\x-3*\x-10)/5});
\draw[thick,domain=-3:9.2] plot (\x,{0.6*\x+1.2});
\draw[ax] (8,0)--(8,6); \draw[ax] (0,6)--(8,6);
\foreach \p in {(-2,0),(5,0),(0,-2),(0,1.2),(8,6)} \fill \p circle (2pt);
\node[above left] at (-2,0){$-2$}; \node[below right] at (5,0){$5$}; \node[left] at (0,-2.1){$-2$};
\node[left] at (0,1.4){$1{,}2$}; \node[below] at (8,0){\color{aux}$8$}; \node[left] at (0,6){\color{aux}$6$};
\node[right] at (8.1,5.6){$(8;6)$};
\end{tikzpicture}}
\textbf{Parabola} $(-2;0)$ va $(5;0)$ dan ötadi: $f(x)=a(x+2)(x-5)$. \ $f(0)=-10a=-2\Rightarrow a=\frac15$:
\[f(x)=\tfrac15(x+2)(x-5).\]
\textbf{Töğri çiziq} $(-2;0)$ va $(0;\,1{,}2)$ dan ötadi: $g(x)=0{,}6x+1{,}2=\tfrac35(x+2)$.

a) Kesişiş: $\tfrac15(x+2)(x-5)=\tfrac35(x+2)\Rightarrow x-5=3\Rightarrow x_0=8$, \ $y_0=0{,}6\cdot8+1{,}2=6$. \ $x_0+y_0=14$.

b) Böyalgan soha $[-2;\,8]$ da $g$ dan pastda, $f$ dan yuqorida:
\[S=\int_{-2}^{8}\big(g(x)-f(x)\big)dx=\frac15\int_{-2}^{8}(x+2)(8-x)\,dx=\frac15\int_{-2}^{8}(-x^2+6x+16)\,dx\]
\[=\frac15\left(-\frac{x^3}3+3x^2+16x\right)\Big|_{-2}^{8}=\frac15\left(\frac{448}{3}+\frac{52}{3}\right)=\frac15\cdot\frac{500}3=\frac{100}3.\]
\Javob{a) 14; \quad b) $\frac{100}3$}

%================= 41
\newpage
\Savol{41}
\fig{\begin{tikzpicture}[scale=0.82]
\coordinate (A) at (0,0); \coordinate (C) at (15*0.85,0); \coordinate (B) at (5.8*0.85,3.919*0.85);
\coordinate (I) at (6*0.85,1.837*0.85); \coordinate (J) at (9*0.85,1.2247*0.85);
\draw[thick] (A)--(B)--(C)--cycle;
\draw[thick] (I) circle (1.837*0.85); \draw[thick] (J) circle (1.2247*0.85);
\draw[ax] (C)--(I);
\draw[axs] (I)--(6*0.85,0); \draw[axs] (J)--(9*0.85,0);
\draw[ax] (6*0.85+0.15,0)--(6*0.85+0.15,0.15)--(6*0.85,0.15);
\draw[ax] (9*0.85+0.15,0)--(9*0.85+0.15,0.15)--(9*0.85,0.15);
\fill (I) circle (1.5pt); \fill (J) circle (1.5pt);
\node[left] at (A){$A$}; \node[above] at (B){$B$}; \node[right] at (C){$C$};
\node[above left] at (I){$I$}; \node[above] at ($(J)+(0.05,0.08)$){$J$};
\node[aux,left] at (6*0.85,0.65){$R$}; \node[aux,right] at (9*0.85,0.5){$r$};
\node[below] at (6*0.85,0){$P$}; \node[below] at (9*0.85,0){$Q$};
\node[above left] at ($(A)!0.5!(B)$){$7$}; \node[above right] at ($(B)!0.5!(C)$){$10$};
\node[below] at (11,0){$15$};
\pic[draw,angle radius=1.6cm,"\scriptsize$\frac C2$",angle eccentricity=1.25] {angle=I--C--A};
\end{tikzpicture}}
a) Yarim perimetr $p=\frac{7+10+15}2=16$. Geron formulasi:
\[S=\sqrt{16\cdot(16-7)(16-10)(16-15)}=\sqrt{16\cdot9\cdot6\cdot1}=12\sqrt6,\qquad R=\frac Sp=\frac{12\sqrt6}{16}=\frac{3\sqrt6}4.\]
b) Ikkala aylana ham $AC$ va $BC$ ga urinadi, demak markazlari $I$ va $J$ — $C$ burçak bissektrisasida (\textcolor{aux}{qizil}).\\
$\cos C=\dfrac{10^2+15^2-7^2}{2\cdot10\cdot15}=\dfrac{276}{300}=\dfrac{23}{25}$, \quad $\sin\dfrac C2=\sqrt{\dfrac{1-\cos C}2}=\sqrt{\dfrac1{25}}=\dfrac15$.\\
$\triangle CPI$ va $\triangle CQJ$ dan: \ $CI=\dfrac{R}{\sin\frac C2}=5R$, \ $CJ=\dfrac{r}{\sin\frac C2}=5r$.\\
Aylanalar taşqi urinadi: $IJ=R+r$, ya'ni $CI-CJ=R+r$:
\[5R-5r=R+r\ \Rightarrow\ 4R=6r\ \Rightarrow\ r=\frac23R=\frac23\cdot\frac{3\sqrt6}4=\frac{\sqrt6}2.\]
\Javob{a) $\frac{3\sqrt6}4$; \quad b) $\frac{\sqrt6}2$}

%================= 42
\newpage
\Savol{42}
\fig{\begin{tikzpicture}[scale=0.85]
\coordinate (A) at (0,0); \coordinate (B) at (9,0); \coordinate (D) at (1.6,4); \coordinate (C) at (6.6,4);
\coordinate (K) at ($(B)!4/9!(C)$); \coordinate (L) at ($(C)!1/3!(B)$);
\fill[msblue!10] (A)--(K)--(L)--(D)--cycle;
\draw[thick] (A)--(B)--(C)--(D)--cycle; \draw[thick] (A)--(K); \draw[thick] (D)--(L);
\coordinate (K1) at ($(A)!(K)!(B)$); \coordinate (L1) at ($(D)!(L)!(C)$);
\draw[ax] (C)--(L1); \draw[ax] (K)--(K1); \draw[ax] (L)--(L1);
\draw[aux] ($(K1)+(-0.18,0)$)--++(0,0.18)--++(0.18,0);
\draw[aux] ($(L1)+(-0.18,0)$)--++(0,-0.18)--++(0.18,0);
\node[below left] at (A){$A$}; \node[below right] at (B){$B$}; \node[above] at (C){$C$}; \node[above left] at (D){$D$};
\node[right] at (K){$K$}; \node[below right] at ($(L)+(0.02,0.05)$){$L$};
\node[aux,left] at ($(K)!0.5!(K1)$){$\frac49h$}; \node[aux,right] at ($(L)!0.5!(L1)$){$\frac h3$};
\node[font=\small] at (6.55,3.35){$3t$}; \node[font=\small] at (7.2,2.25){$2t$}; \node[right,font=\small] at ($(K)!0.5!(B)+(0.1,0)$){$4t$};
\node[below] at (4.5,0){$9m$}; \node[above] at (4.1,4){$5m$};
\end{tikzpicture}}
$\frac{CL}{LK}=\frac32$, $\frac{KB}{LK}=2$ \ $\Rightarrow$ \ $CL:LK:KB=3:2:4$, ya'ni $BC=9t$.\\
$DC=5m$, $AB=9m$, balandlik $h$: \ $S_{ABCD}=\dfrac{5m+9m}2\cdot h=7mh=42\Rightarrow mh=6$.

a) $K$ dan $AB$ gaça masofa: $BK=\frac49BC$ bölgani uçun $\frac49h$ (\textcolor{aux}{qizil} balandlik):
\[S_{AKB}=\frac12\cdot9m\cdot\frac49h=2mh=12.\]
b) $L$ dan $DC$ gaça masofa: $CL=\frac39BC$ bölgani uçun $\frac39h=\frac h3$:
\[S_{DLC}=\frac12\cdot5m\cdot\frac h3=\frac56mh=5.\]
\[S_{AKLD}=S_{ABCD}-S_{AKB}-S_{DLC}=42-12-5=25.\]
\Javob{a) 12; \quad b) 25}

%================= 43
\newpage
\Savol{43}
\fig{\begin{tikzpicture}[scale=0.95]
\def\R{3.536}
\coordinate (A) at (0,0); \coordinate (B) at (\R,0); \coordinate (C) at (0,\R); \coordinate (D) at (0.990,3.394);
\draw[thick] (A)--(B) arc (0:90:\R)--cycle;
\draw[thick] (C)--(D)--(B); \draw[thick] (A)--(D);
\draw[ax] (C)--(B);
\draw (0.18,0)--(0.18,0.18)--(0,0.18);
\pic[draw=aux,thick,angle radius=0.35cm] {angle=C--D--B};
\pic[draw,angle radius=0.9cm,"\scriptsize$\alpha$",angle eccentricity=1.25] {angle=D--A--C};
\node[left] at (A){$A$}; \node[below] at (B){$B$}; \node[above left] at (C){$C$}; \node[above right] at (D){$D$};
\node[above] at ($(C)!0.5!(D)+(0,0.05)$){$1$}; \node[right] at ($(D)!0.5!(B)+(0.05,0.05)$){$3\sqrt2$};
\node[aux,below left] at ($(C)!0.55!(B)$){$5$};
\end{tikzpicture}}
$\angle CAB=90^\circ$ — markaziy burçak. $D$ nuqta $CB$ yoyida, şuning uçun $\angle CDB$ — katta yoyga ($270^\circ$) tiralgan içki çizilgan burçak:
\[\angle CDB=\tfrac12\cdot270^\circ=135^\circ.\]
$\triangle CDB$ da kosinuslar teoremasi (\textcolor{aux}{qizil} $CB$ vatar):
\[CB^2=1^2+(3\sqrt2)^2-2\cdot1\cdot3\sqrt2\cdot\cos135^\circ=1+18+6=25\ \Rightarrow\ CB=5.\]
$\triangle CAB$ — teng yonli töğri burçakli: $R\sqrt2=5\Rightarrow R=\frac5{\sqrt2}$, \ $R^2=\frac{25}{2}$.

a) $\triangle CAD$ da $AC=AD=R$, $CD=1$, $\alpha=\angle CAD$. Kosinuslar teoremasi:
\[1=R^2+R^2-2R^2\cos\alpha\ \Rightarrow\ \cos\alpha=1-\frac1{2R^2}=1-\frac1{25}=\frac{24}{25},\qquad \sin\alpha=\sqrt{1-\frac{576}{625}}=\frac7{25}.\]
b) \[S_{CAD}=\frac12R^2\sin\alpha=\frac12\cdot\frac{25}2\cdot\frac7{25}=\frac74.\]
\Javob{a) $\frac7{25}$; \quad b) $\frac74$}

%================= 44
\newpage
\Savol{44}
\fig{\begin{tikzpicture}[scale=0.62]
\def\r{3} \def\e{0.9}
\draw[thick] (-\r,0)--(-\r,6); \draw[thick] (\r,0)--(\r,6);
\draw[thick] (0,6) ellipse ({\r} and {\e});
\draw[thick] (-\r,0) arc (180:360:{\r} and {\e}); \draw[dashed] (\r,0) arc (0:180:{\r} and {\e});
\draw[ax] (\r,3) arc (0:180:{\r} and {\e}); \draw[axs] (-\r,3) arc (180:360:{\r} and {\e});
\coordinate (T) at (0,6); \coordinate (Bo) at (0,0);
\coordinate (P1) at (-3,3); \coordinate (P2) at (3,3); \coordinate (P3) at (0,{3-\e}); \coordinate (P4) at (0,{3+\e});
\draw[thick,okgreen] (P1)--(P3)--(P2); \draw[dashed,okgreen] (P2)--(P4)--(P1);
\foreach \p in {P1,P2,P3} {\draw[thick,okgreen] (T)--(\p); \draw[thick,okgreen] (Bo)--(\p);}
\draw[dashed,okgreen] (T)--(P4); \draw[dashed,okgreen] (Bo)--(P4);
\draw[ax] (P1)--(P2); \draw[ax] (T)--(Bo);
\fill (T) circle (2pt); \fill (Bo) circle (2pt); \fill (0,3) circle (1.5pt);
\node[aux,right] at (3.05,3){$a\sqrt2$}; \node[aux,right] at (0.05,4.6){$h$};
\node[okgreen] at (-1.5,2.2){$a$};
\end{tikzpicture}}
\textcolor{okgreen}{Oktaedr} ikki uçi silindr asoslari markazlarida, qolgan 4 uçi esa örta kesim aylanasidagi kvadratning uçlarida yotadi.

a) Muntazam oktaedr — asoslari umumiy ikkita muntazam törtburçakli piramida. Kvadrat diagonali $a\sqrt2$, har bir piramida balandligi $\frac{a\sqrt2}{2}$:
\[V=2\cdot\frac13a^2\cdot\frac{a\sqrt2}2=\frac{a^3\sqrt2}{3}=36\ \Rightarrow\ a^3=54\sqrt2=(3\sqrt2)^3\ \Rightarrow\ a=3\sqrt2.\]
b) Silindr radiusi = kvadratga taşqi çizilgan aylana radiusi: $r=\frac{a\sqrt2}2=3$.\\
Silindr balandligi = oktaedrning qarama-qarşi uçlari orasidagi masofa: $h=a\sqrt2=6$.
\[V_{\text{sil}}=\pi r^2h=\pi\cdot9\cdot6=54\pi.\]
\Javob{a) $3\sqrt2$; \quad b) $54\pi$}

%================= 45
\newpage
\Savol{45}
\fig{\begin{tikzpicture}[scale=0.9]
\fill[gray!35] (0,0) rectangle (4,2);
\fill[white] (1,1) circle (1) (3,1) circle (1);
\draw[thick,fill=msblue!15] (1,1) circle (1) (3,1) circle (1);
\fill[msblue!25] (4,0) rectangle (10,2);
\draw[thick] (0,0) rectangle (10,2);
\draw[axs] (2,0)--(2,2); \draw[axs] (4,0)--(4,2);
\draw[ax] (1,1)--(2,1); \node[aux,above] at (1.5,1){\small$r$};
\node at (7,1){yon sirt $2\pi r\times h$};
\draw[|<->|] (0,-0.45)--node[lbl]{$4$}(4,-0.45); \draw[|<->|] (4,-0.45)--node[lbl]{$6$}(10,-0.45);
\node[left] at (0,1){$2$};
\end{tikzpicture}}
Asos doiralari radiusi $r$ — ular eni 2 bölgan listga siğişi kerak: $2r\le2\Rightarrow r\le1$. Yon sirt — ölçamlari $2\pi r\times h$ bölgan töğri törtburçak.

Çiqindi eng kam bölişi uçun $r=1$ olamiz: ikkita doira $2\times2$ li kvadratlarga joylaşib, $4\times2$ qismni egallaydi. Qolgan $6\times2$ qism yon sirt böladi: $2\pi r=2\cdot3\cdot1=6$ — aynan mos, $h=2$.

a) \[V=\pi r^2h=3\cdot1^2\cdot2=6.\]
b) Metall yuzi $2\cdot10=20$. Işlatilgan qism: $2\pi r^2+2\pi rh=2\cdot3\cdot1+6\cdot2=6+12=18$.\\
Çiqindi (\textcolor{gray}{kulrang} burçaklar): $20-18=2$, ya'ni $\dfrac{2}{20}\cdot100\%=10\%$.
\Javob{a) 6; \quad b) 10\%}

\ShowJavoblar
\end{document}
